\begin{scanpage}{0115}{98}
\[\Sigma^0=\operatorname{tg}\omega\,e^{2ikl}.\tag{17.12}\]
Thus even in the case of short waves ($ka\gg\operatorname{tg}^2\omega$, but $R\gg kl^2$), our Expression (17.08) does not change into the physical optics equation, but substantially differs from it because
\[\sigma=\pi a^2\left|\frac{\frac2n\sin\frac\pi n}{\cos\frac\pi n-\cos\frac{2\omega}n}\right|^2\tag{17.13}\]
and
\[\sigma^0=\pi a^2\operatorname{tg}^2\omega.\tag{17.14}\]
With this
\[\sigma=\sigma^0\left|\frac{\frac2n\sin\frac\pi n}{(\cos\frac\pi n-\cos\frac{2\omega}n)\operatorname{tg}\omega}\right|^2,\tag{17.15}\]
that is, for sufficiently short waves (or for sufficiently large dimensions of the cone) the function $\sigma$ is proportional to $\sigma^0$. The coefficient of proportionality here does not depend on the cone dimensions, but is determined only by its shape.

This result is graphically illustrated by the curves giving the effective scattering area of a cone ($\omega=10^\circ25'$, $k=\pi$, $\Omega=90^\circ$) as a function of its length (Figure 31). Whereas our equation (the continuous line) is in satisfactory agreement with the results of measurements (the small crosses)\textsuperscript{(1)}, the physical optics approach (the dashed line) gives values which are smaller than the experimental values by 13--15 dB. For sharply pointed cones, the nonuniform part of the current has an especially large values. In Figure 32, a curve is constructed for the effective surface of a cone ($ka=2.75\pi$, $\Omega=90^\circ$)

\vfill\noindent\rule{.5\linewidth}{.4pt}

Footnote (1) appears on page 113.
\end{scanpage}
